\documentclass[10pt,twocolumn]{article}

\usepackage[margin=0.72in,columnsep=0.25in]{geometry}
\usepackage{amsmath,amssymb}
\usepackage{booktabs}
\usepackage{graphicx}
\usepackage{microtype}
\usepackage{natbib}
\usepackage{tabularx}
\usepackage{xcolor}
\usepackage{enumitem}
\usepackage{xurl}
\usepackage[hidelinks]{hyperref}
\usepackage[T1]{fontenc}
\usepackage{lmodern}
\usepackage{titlesec}

\hypersetup{
  pdftitle={ProofPathBench: A Benchmark for Verifiability-Aware Planning by Tool-Using Language Agents},
  pdfauthor={Vamshi Krishna Madhavan},
  pdfsubject={Benchmark design and deterministic validation for verification-aware agent planning},
  pdfkeywords={language agents, tool use, verification, benchmark, reliability}
}

\definecolor{protocolblue}{HTML}{1769AA}
\definecolor{protocolgray}{HTML}{475569}
\setlist[itemize,enumerate]{nosep,leftmargin=*}
\setlength{\parskip}{2.5pt}
\setlength{\parindent}{1em}
\widowpenalty=10000
\clubpenalty=10000
\displaywidowpenalty=10000
\renewcommand{\topfraction}{0.9}
\renewcommand{\bottomfraction}{0.8}
\renewcommand{\textfraction}{0.08}
\renewcommand{\floatpagefraction}{0.8}
\titleformat{\section}{\large\bfseries}{\thesection}{0.65em}{}
\titleformat{name=\section,numberless}{\large\bfseries}{}{0pt}{}
\titleformat{\subsection}{\normalsize\bfseries}{\thesubsection}{0.65em}{}
\titlespacing*{\section}{0pt}{1.35ex plus 0.5ex minus 0.2ex}{0.75ex}
\titlespacing*{\subsection}{0pt}{1.0ex plus 0.4ex minus 0.2ex}{0.5ex}

\newcommand{\system}{\textsc{ProofPath}}
\newcommand{\bench}{\textsc{ProofPathBench}}
\newcommand{\vpsr}{\textsc{VPSR}}
\newcommand{\vts}{\textsc{VTS}}
\newcommand{\fcr}{\textsc{FCR}}

\title{\vspace{-1.1em}\resizebox{0.98\textwidth}{!}{\textbf{ProofPathBench: A Benchmark for Verifiability-Aware Planning by Tool-Using Language Agents}}}
\author{
  Vamshi Krishna Madhavan\\
  Independent Researcher\\
  \href{mailto:vamshi-madhavan@outlook.com}{\texttt{vamshi-madhavan@outlook.com}}\\
  \href{https://vamshi0104.github.io/proofpathbench/}{\texttt{vamshi0104.github.io/proofpathbench/}}
}
\date{Benchmark release: October 2026}

\begin{document}
\maketitle

\begin{abstract}
Tool-using language agents can change files, accounts, and other persistent state. A tool
may report success even when the intended change did not occur. We introduce \bench, an
implemented benchmark that tests whether an agent chooses a plan with stronger
independent evidence before execution begins. This voluntary, ex-ante choice is distinct
from detecting or recovering from a failure after a tool has run. The benchmark contains 96 synthetic
scenarios across eight domains. Each scenario offers two plans for the same goal and
varies the simulated cost and delay of verification, the stated failure risk, the
consequence severity, and the instruction given to the agent. The environment keeps the
true state separate from tool responses and from an evaluator-only trace. Automated
checks cover 384 task-treatment units, 768 no-fault plan executions, state isolation,
balanced presentation, and nine seeded failure types. A prospective power analysis also
defines a 2,304-choice reference evaluation and shows which comparisons it can support.
This paper presents the benchmark, its outcome definitions, its deterministic validation,
and its evaluation design. The complete scenarios, implementation, validation records,
and release checks accompany the paper. It does not report results from language models
and makes no claim about how current agents value verification.
\end{abstract}

\section{Introduction}

Language-model agents increasingly use tools to modify files, databases, profiles,
calendars, messages, and configurations. This creates a basic reliability problem: a
tool's reply may not match the state of the outside world. A write can report success
after doing nothing, time out after committing, update only part of an object, or change
the wrong target. If an agent treats a plausible reply as proof, it may confidently claim
that an unfinished task is complete.

Prior work studies stateful tool-use evaluation, runtime failures, false success,
postcondition checking, transactional execution, and verification-aware interfaces
\citep{trivedi2024appworld,yao2025taubench,lu2025toolsandbox,sun2024tools,
advani2026confident,wang2026callability,mansoor2026verified,liu2026toolgate}.
We ask a narrower question. When two plans pursue the same goal, will an agent choose the
plan that makes the result easier to verify independently? How much simulated cost or
delay will it accept for that evidence?

\bench{} turns this question into a fixed choice made before execution. The benchmark is
fully synthetic: it makes no real purchase, sends no message, and changes no real account
or infrastructure. Its main rule is simple: a tool response is not ground truth. The
evaluator checks the final state separately from everything shown to the agent.

This work contributes:

\begin{itemize}
  \item a task-specific definition of outcome verifiability with six dimensions;
  \item 96 typed, deterministic scenarios with paired plans and counterbalanced
        presentations;
  \item a controlled design for measuring how verification choices change with price,
        stated risk, and consequence severity;
  \item an implemented mock environment, fault injector, evaluator, immutable run logger,
        and analysis pipeline; and
  \item deterministic validation and a power analysis that show what the reference
        evaluation can and cannot support.
\end{itemize}

The contribution is the benchmark and its validated research artifact. We do not claim
that any agent or model family already prefers verifiable plans. Model comparisons and a
possible \system{} planner are outside the scope of this release.

\section{Problem Formulation}

Let a task be
\begin{equation}
T=(s_0,g,\phi_g,\kappa,B),
\end{equation}
where $s_0$ is authoritative initial state, $g$ is the semantic goal,
$\phi_g(s)$ is the authoritative postcondition, $\kappa$ is a set of collateral
constraints, and $B$ is an execution budget. A plan $P$ is a sequence of mutation,
read-only, and control operations with declared simulated cost and latency.

Two plans count as matched only if a deterministic no-fault run shows that both satisfy
$\phi_g$ and $\kappa$, use the same authorization and inputs, and differ only in their
evidence steps and stated overhead. An extra read or status check cannot change the
target or value of the write.

\subsection{Outcome verifiability}

An evidence item $e$ is evaluated relative to $T$ using
\begin{equation}
Q(e,T)=(I,A,S,F,L,C),
\end{equation}
Here, $I$ measures independence from the write response and its failure boundary; $A$
measures source authority; $S$ measures how directly the evidence addresses the desired
state; $F$ measures freshness; $L$ links the evidence to the intended entity or action;
and $C$ measures coverage of the goal and its constraints. Each dimension is rated
\emph{none}, \emph{weak}, or \emph{strong}. Every rating has a written reason and a
provenance graph. We do not combine the dimensions with an arbitrary weighted score.

Plan $P_2$ evidence-dominates $P_1$ when its best valid evidence is at least as strong on
all six dimensions and stronger on at least one. Every primary item must have one clear
dominant plan. A write response such as \texttt{\{"status":"success"\}} has
$I=\mathrm{none}$, so it cannot provide independent evidence by itself.

\subsection{Estimands}

The main outcome is the Verifiable Plan Selection Rate (\vpsr): the share of eligible
choices that select the evidence-dominant plan. The main comparisons are the baseline
selection rate, the effect of an explicit verification instruction, and the effect of a
higher verification premium. The environment also supports
Verified Task Success (\vts), False Completion Rate (\fcr), recovery success, uncertainty
fidelity, and cost and delay overhead. These measures require model runs and are defined
here for future evaluations.

\begin{figure*}[t]
  \centering
  \includegraphics[width=0.94\textwidth]{figures/proofpath_architecture.pdf}
  \caption{Measurement architecture. Plan registration occurs before execution. The
  tool-visible response may diverge from authoritative state; the evaluator scores the
  latter using an evaluator-only oracle trace. The diagram is schematic, not an empirical
  result.}
  \label{fig:architecture}
\end{figure*}

\section{Related Work}

\paragraph{Stateful tool-use benchmarks.}
AppWorld evaluates cross-application tasks against state-based unit tests and collateral
constraints \citep{trivedi2024appworld}. \(\tau\)-bench measures tool-agent-user
interaction over persistent domain state \citep{yao2025taubench}, while ToolSandbox
evaluates conversational tool use through state milestones \citep{lu2025toolsandbox}.
ToolEmu uses model-emulated sandboxes to expose risky agent behavior
\citep{ruan2024toolemu}. Proxy State-Based Evaluation instead infers structured state
from interaction traces and validates it with model-based judges
\citep{chuang2026proxy}. Together, these studies show why final state should be checked
directly. They do not isolate a voluntary choice between matched plans that mainly differ
in the evidence they produce.

\paragraph{Tool failures and false completion.}
Tools Fail studies whether models detect silent errors in faulty tools
\citep{sun2024tools}; more recent work examines dynamic replanning, production API failure,
recoverable tool-environment hazards, and confident false-success claims
\citep{zhu2026when,tian2026toolbenchx,li2026silentprobe,advani2026confident}. This line
asks whether agents detect, recover from, or accurately report failures. \bench{} instead
measures demand for evidence before the agent knows whether a failure occurred.

\paragraph{Selection and presentation effects.}
ToolTweak shows that names and descriptions can strongly bias selection among
functionally similar tools \citep{sneh2025tooltweak}. This result motivates the neutral
aliases, reversed plan order, prompt paraphrases, and presentation covariates used here.
Those controls reduce obvious presentation bias, but they do not prove that every
remaining wording effect has been removed.

\paragraph{Verification and execution mechanisms.}
Controlled interface interventions, verified tool-call wrappers, contract-grounded
execution, and semantic enterprise APIs can improve operational reliability
\citep{wang2026callability,mansoor2026verified,liu2026toolgate,pan2026agentfirst}.
Pre-execution verifiers for robot plans and shell commands audit or revise a proposed
action rather than ask whether an agent voluntarily pays for outcome evidence
\citep{grigorev2025verifyllm,liu2026care}. Recent work also separates planning-mode
selection from faithful execution, but studies planning structure rather than
outcome-evidence demand \citep{oota2026declaration}.

Transactional systems such as Atomix and Cordon address settlement, isolation, and
validation of tool effects \citep{mohammadi2026atomix,chen2026cordon}. AgentOps tracks
the agent lifecycle \citep{dong2024agentops}. These systems are close technical
precedents. In most cases, however, the system assigns the safeguard; the agent does not
choose it at a controlled range of prices.

Within the sources checked, we found no work that combines matched plan choices before
execution, an explicit verification premium, risk and severity factors, and scoring from
authoritative state. Tool-selection bias, unreliable-tool recovery, postcondition-aware
runtimes, and state-based evaluation are established neighboring areas rather than claims
of novelty here. The narrower conclusion is limited to our structured search and may
change as new work appears.

\section{ProofPathBench}

\subsection{Scenario composition}

The benchmark contains 96 base scenarios, with 12 in each of eight domains. Every
operation runs in a local mock environment. Severity is represented through concrete
simulated effects, such as dependents, attendees, recipients, service instances,
fulfillment penalties, or entitlement periods. It is not represented by a label alone.

\begin{table}[t]
  \centering
  \caption{Benchmark composition generated from scenario manifests.}
  \label{tab:composition}
  \begin{tabular}{lr}
\toprule
Domain & Base scenarios \\
\midrule
Calendar & 12 \\
Commerce & 12 \\
Configuration & 12 \\
Database & 12 \\
Filesystem & 12 \\
Messaging & 12 \\
Profile & 12 \\
Subscription & 12 \\
\midrule
Total & 96 \\
\bottomrule
\end{tabular}

\end{table}

Each item defines a goal, the initial state, success and safety conditions, two candidate
plans, evidence profiles, failure rules, cost and delay, and four balanced presentation
versions. Neutral plan names and plan order rotate so that the stronger option is not
tied to one label or position.

The stronger plan uses an independent readback, an operation-status lookup, or evidence
from more than one source. Each method appears in 32 base scenarios. The analysis adjusts
for the method but does not make causal claims about one method versus another.

\subsection{Evaluation design}

Risk and severity are balanced across the 96 base tasks. Each task appears at all four
verification premiums. This produces 384 task-treatment units and 16 units in every
premium-by-risk-by-severity cell. The four presentation versions rotate across premiums
within each task. Cost and delay change together, so the design measures one combined
execution premium rather than separate cost and delay effects.

\begin{table}[t]
  \centering
  \caption{Reference evaluation factors.}
  \label{tab:factors}
  \begin{tabularx}{\columnwidth}{@{}lX@{}}
\toprule
Factor & Levels / role \\
\midrule
Premium & 1.00, 1.10, 1.50, 2.00; within task \\
Disclosed risk & 0, 0.05, 0.20; between task \\
Severity & Low, high; between task \\
Instruction & Vanilla, verify instruction \\
Evidence & Readback, status lookup, multiple source; nuisance \\
Repetitions & 1 per unit \\
\bottomrule
\end{tabularx}

\end{table}

\subsection{Failure model and state separation}

The environment implements nine seeded faults: explicit failure, false success, timeout
before execution, timeout after execution, partial update, stale readback, delayed
visibility, wrong-target update, and duplicated action. Each seed is derived from the
master seed and fixed run identifiers. The environment keeps three state layers separate:

\begin{enumerate}
  \item authoritative state, used for outcome scoring;
  \item tool-visible state and responses, which may be faulty; and
  \item an append-only evaluator-only oracle trace.
\end{enumerate}

This design can return a valid-looking success message even when the goal is false. It
can also return a timeout after the write has already committed.

\section{Deterministic Validation}

The validation suite checks schemas, content hashes, factor balance, hidden label
leakage, name and order rotation, evidence dominance, cost, and delay. Runtime tests run
both plans without faults in all 384 treatment units. They confirm that both plans reach
the same goal and satisfy the same safety conditions. Other tests confirm that tool
responses cannot change evaluator state, every fault activates under a known seed, and
both passing and failing conditions are exercised.

\begin{table}[t]
  \centering
  \caption{Deterministic acceptance evidence generated from the runtime report.}
  \label{tab:validation}
  \begin{tabular}{lr}
\toprule
Deterministic check & Value \\
\midrule
Generated base scenarios & 96 \\
Split-plot treatment units & 384 \\
No-fault plan executions & 768 \\
Injectable failure classes & 9 \\
Runtime validation errors & 0 \\
\bottomrule
\end{tabular}

\end{table}

At the release freeze, all 48 automated tests pass, along with static linting and strict
type checking. The runtime report has no validation errors. It covers 768 no-fault plan
executions and activates all nine failure types. These results validate the software and
the benchmark design. They are not observations of language-model behavior and do not
test a behavioral hypothesis.

Mechanical checks show that the implementation is internally consistent. They cannot
establish full construct validity. Before interpreting model scores, researchers should
still use independent human review to check task equivalence, evidence strength, and
wording bias. The release includes a fixed, blinded 32-item protocol for two independent
reviewers, balanced at four items per domain. It prespecifies completeness rules,
agreement statistics, thresholds, and adjudication. The review has not been run, so no
human-validation result is claimed. We therefore describe this release as implemented
and mechanically validated.

\section{Evaluation Framework}

\subsection{Pre-execution registration}

In an evaluation run, the agent receives a task, two neutrally named plans, their
simulated cost and delay, and the stated risk. It must choose one plan before execution.
The choice cannot be changed. Model mistakes remain part of the results. A run is excluded
only for a provider failure before any response, invalid structured output after the
fixed repair allowance, or a recorded infrastructure failure.

The reference evaluation contains
\[
96\ \text{tasks}\times4\ \text{premiums}\times2\ \text{instructions}
\times3\ \text{models}=2{,}304
\]
registered choices. The system stores exact model versions, parameters, prompts, schemas,
seeds, raw responses, and exclusions without overwriting earlier records. This release
tests the pipeline with fixtures but does not call provider models.

\subsection{Hypotheses and analysis}

The pilot defines two confirmatory tests in advance:

\begin{description}
  \item[H1:] vanilla \vpsr{} at zero premium exceeds chance;
  \item[H2:] an explicit verification instruction increases \vpsr{}.
\end{description}

The premium slope (H3), premium-by-risk term (H4), and premium-by-severity term (H5) are
directional estimation targets, not confirmatory pilot tests. H1 is the marginal vanilla,
zero-premium rate contrast against 0.5. The other estimates use logistic regression with
premium, risk, severity, instruction, and the two planned interactions, adjusting for
model, domain, presentation, evidence method, and plan order. H3 is reported as the
design-average log-odds slope rather than a single reference-cell coefficient. The
three-model pilot uses model fixed effects and scenario-cluster-robust uncertainty; three
model clusters are too few for a stable model-cluster sandwich estimate. A later,
broader study may use two-way clustering when the number of distinct model clusters is
adequate. Holm correction covers H1 and H2 only. Every result must report raw rates,
denominators, exclusions, effect sizes, and confidence intervals.

The smallest effect of interest is fixed at 10 percentage points, and the analysis formula
is set before model data are collected. Comparisons by model or evidence method, unusual
curve shapes, rationales, and open-action traces remain exploratory unless a separate
study has enough power for them.

\section{Evaluation Sizing and Sensitivity}

The power analysis does not use language-model outcomes. It uses 10,000 Monte Carlo runs
for the instruction comparison and a scenario-clustered linear-probability approximation
for the factorial terms. It assumes two confirmatory one-sided tests with Holm correction
and the documented random-effect and Bernoulli structure. Power for H1 is not simulated:
its precision depends on the as-yet unobserved baseline rate and cluster variation.

\begin{figure}[t]
  \centering
  \includegraphics[width=\columnwidth]{figures/prospective_power.pdf}
  \caption{Prospective power under stated simulation assumptions. These curves are not
  empirical model results. Repetitions increase call-level precision but do not add new
  task archetypes.}
  \label{fig:power}
\end{figure}

\begin{table*}[t]
  \centering
  \caption{Prospective estimated power. ``Instr.'' is the instruction contrast; the
  final column reports premium-by-risk and premium-by-severity interaction estimates.
  Values are simulation outputs, not observed effects.}
  \label{tab:power}
  \begin{tabular}{rrrrr}
\toprule
Rep. & Calls & Instr. & Premium & Risk / severity int. \\
\midrule
1 & 2,304 & 0.991 & 0.484 & 0.290 / 0.377 \\
2 & 4,608 & 1.000 & 0.769 & 0.501 / 0.634 \\
3 & 6,912 & 1.000 & 0.909 & 0.670 / 0.802 \\
\bottomrule
\end{tabular}

\end{table*}

With one repetition, estimated power is 0.991 for the assumed instruction effect. Power
is approximately 0.484 for premium, 0.290 for the premium-by-risk interaction, and 0.377 for
the premium-by-severity interaction under the stated assumptions. The exact generated
values appear in Table~\ref{tab:power}. The 2,304-choice design is therefore useful for
the instruction comparison and for estimating benchmark variance, but not for definitive
premium or interaction claims. A later confirmatory study should use blinded pilot
estimates to set its sample size.

\section{Artifact Controls}

The arXiv ancillary archive contains the generated scenarios, JSON schema, hash index,
configuration snapshots, seed rules, mock environments, fault schedules, evaluation
conditions, run-directory checks, analysis code, tests, bibliography audit, and
related-work matrix. The script \texttt{paper/build\_assets.py} builds the paper's figures
and tables from stored JSON and YAML files.

The accompanying confound report records operation count, serialized plan length,
distinct tools, extra verification operations, cost, latency, aliases, plan order, and
evidence mechanism for every scenario. Alias set, prompt template, evidence method, and
dominant-plan display position are balanced. The diagnostic also confirms two residual
couplings: the evidence-dominant plan is necessarily longer, and simulated cost and
latency are perfectly correlated. Neither effect can be identified separately in this
version.

Evaluation commands reject incomplete preregistration and review files. Fixture outputs
are labeled non-empirical and cannot enter the confirmatory analysis. Provider responses
go into fixed, uniquely named run directories. Versioned artifacts exclude secrets and
API keys. These controls reduce the chance of mixing fixtures, simulations, and model
observations. They do not replace independent replication or scientific review.

The ancillary archive is self-contained, carries the Apache License 2.0, and includes
the exact scenario manifests and validation reports used by this paper. The release
build recreates figures and tables from machine-readable records, compiles the manuscript
from a clean unpacked source archive, compares extracted text and page count with the
local build, and rejects unresolved citations, missing fonts, unsafe filenames,
placeholders, and private or generated cache files.

\section{Limitations}

\paragraph{Scope of the current release.}
This release provides evidence about benchmark construction and deterministic execution.
It does not include a provider-model evaluation. It therefore cannot show whether agents
prefer verifiable plans, whether instructions change that preference, or whether
verification improves task success.

\paragraph{Construct validation.}
Mechanical equivalence does not rule out every wording, salience, or perceived-capability
bias. A behavioral study should complete the released blinded human review and a
prospectively registered manipulation check before interpreting differences between
conditions.

\paragraph{External review.}
The current release has not undergone peer review or the planned blinded two-reviewer
assessment of semantic equivalence and evidence dominance. Automated tests cannot replace
those judgments. This paper should therefore be read as a benchmark-and-methods release,
not as validated evidence about model behavior.

\paragraph{Synthetic incentives.}
Cost, delay, failure probability, and consequences are simulated. A language model does
not bear these costs in the same way as a person or organization. Its choice may reflect
prompt associations rather than a stable attitude toward risk.

\paragraph{Joint premium.}
Cost and delay change together, so this study cannot separate their effects. That would
require a design that varies them independently. Plan length, tool familiarity, and the
extra operation may also influence the choice despite neutral names and balanced order.

\paragraph{Evidence independence.}
Evidence independence depends on the task and system design. A read endpoint backed by
the same faulty subsystem may add little useful confirmation. The categorical rubric
avoids false numerical precision, but it still includes author judgment.

\paragraph{Generalization and drift.}
Ninety-six synthetic tasks cannot support broad claims about real systems, agent designs,
or model families. Provider changes and nondeterministic behavior may also make later
replication harder, even with fixed seeds.

\paragraph{Literature coverage.}
The related-work audit is a structured rapid review, not a complete systematic review.
Several closely related sources are recent preprints. New work may change the novelty
boundary.

\section{Ethics}

All benchmark actions run in deterministic mock environments. The package makes no real
payment, sends no real message, and changes no real account or infrastructure. Scenarios
use synthetic entities and limited simulated consequences. The main risk is
misinterpretation: readers could mistake a polished benchmark for completed behavioral
evidence. The abstract, validation section, and conclusion therefore state clearly what
was and was not tested.

A future model evaluation would have financial and environmental costs. It should use a
declared budget, bounded outputs, and the smallest design that answers the prospectively specified
question. Negative and null results should be retained. This benchmark alone is not
enough evidence to deploy a system as a reliability safeguard.

\section{Conclusion}

\bench{} is an implemented benchmark for studying whether tool-using language agents
choose plans that make external outcomes easier to verify independently. Before
execution, the agent chooses between two matched plans. The benchmark varies the price of
evidence, the stated failure risk, and the consequence severity. It keeps the true state
separate from the response shown to the agent. The released software passes deterministic
validation across 384 task-treatment units, 768 no-fault plan executions, and nine
injected failure types. These results show that the benchmark works as specified. They do
not answer the behavioral question. The artifact supports transparent, comparable model
evaluations while keeping simulated and prospective values separate from observed
results.

\section*{AI Assistance Disclosure}

Generative-AI tools assisted with software development, literature-search support,
manuscript drafting and language editing, and figure layout. No AI system is listed as an
author. The named human author remains responsible for the research design,
source verification, claims, code, and final text.

\section*{Author Contributions}

The author designed the study, implemented and validated the software and benchmark,
conducted the literature audit and prospective analysis, and prepared the manuscript and
release artifact.

\section*{Data and Code Availability}

The complete benchmark, source code, configuration files, deterministic validation
records, tests, and analysis scaffolding are included as an arXiv ancillary archive. The
software is released under the Apache License 2.0. No proprietary model output, personal
data, credentials, or undisclosed empirical results are included.

\bibliographystyle{plainnat}
\begingroup
\small
\bibliography{references}
\endgroup

\appendix
\clearpage
\onecolumn

\section{Failure Semantics}

\begin{table}[h]
\centering
\caption{Implemented failure classes and their intended evaluator semantics.}
\begin{tabularx}{\textwidth}{@{}lXX@{}}
\toprule
Failure & Authoritative transition & Agent-visible behavior \\
\midrule
Explicit failure & No intended change & Explicit error \\
False success & No intended change & Syntactically valid success \\
Timeout before & No intended change & Timeout or unknown status \\
Timeout after & Intended change commits & Same timeout or unknown status \\
Partial update & Subset of fields commits & Success, partial, or timeout by condition \\
Stale readback & Current state exists & Older version is returned \\
Delayed visibility & Write commits; read model lags & Old value until the seeded visibility step \\
Wrong target & Unintended entity changes & Plausible success for the requested call \\
Duplicated action & Logical effect applies twice & One or ambiguous response \\
\bottomrule
\end{tabularx}
\end{table}

\section{Outcome Definitions}

\begin{description}
  \item[VPSR:] evidence-dominant registrations divided by eligible registered choices.
  \item[VTS:] authoritative postcondition and all collateral constraints hold at termination.
  \item[FCR:] the agent asserts completion while \vts{} is false.
  \item[Recovery success:] after an activated fault, the goal is satisfied without a
        prohibited duplicate or collateral effect.
  \item[Uncertainty fidelity:] the terminal claim distinguishes verified success,
        verified failure, and unresolved status correctly.
\end{description}

Every future metric report must include its numerator, denominator, exclusions,
confidence interval, and scenario/model clustering strategy.

\par\vspace{1.2ex}
\section{Artifact Map}

The source package is backed by the following repository paths:

\begin{itemize}
  \item \texttt{benchmark/scenarios/}: 96 generated task manifests;
  \item \texttt{benchmark/runtime\_validation\_report.json}: deterministic runtime evidence;
  \item \texttt{configs/pilot.yaml}: proposed empirical configuration and gates;
  \item \texttt{proofpath/environments/} and \texttt{proofpath/failures/}: mock state and faults;
  \item \texttt{proofpath/evaluation/}: outcomes, metrics, and inferential analysis;
  \item \texttt{research/power\_analysis.json}: prospective power outputs;
  \item \texttt{research/confound\_diagnostics.json}: design-coupling diagnostics;
  \item \texttt{research/human\_validation/}: blinded two-reviewer protocol and blank form;
  \item \texttt{research/related\_work\_matrix.md}: structured comparison; and
  \item \texttt{research/citation\_audit.md}: bibliography verification status.
\end{itemize}

\par\vspace{1.2ex}
\section{Release Scope Statement}

The release contains fixture smoke tests but no OpenAI or other provider-model evaluation.
Fixtures are explicitly marked non-empirical. The manuscript includes no blank,
simulated, or prospective quantity as if it were an observed behavioral result.

\end{document}
